% =========================================================================
% 附录：支撑材料清单、软件环境、图表对应关系与核心程序
% =========================================================================

\section{支撑材料文件列表}

\subsection{输入数据（题面附件）}

\begin{itemize}
  \item \texttt{data/附件1.xlsx}：代表日逐 10 分钟的电价、小区负载与光伏预测功率（题面附件 1）；
  \item \texttt{data/附件2.xlsx}：2025 年逐日实测负载（工作表 1）与光伏功率（工作表 2）（题面附件 2）；
  \item \texttt{data/附件3.xlsx}：0:00、6:00、12:00、18:00 发布的未来 24 小时整点光伏功率预报（题面附件 3）；
  \item \texttt{data/附件4.xlsx}：2025 年逐日逐 10 分钟区间的实时电价（题面附件 4）；
  \item \texttt{data/附件5/}：官方提交模板（\texttt{result1.xlsx}、\texttt{result2.xlsx}、\texttt{result3.xlsx}、\texttt{result4-2.xlsx}、\texttt{result4-3.xlsx}），用于核对各工作簿的工作表名与原始行序。原始数据与模板均按题目要求随支撑材料提交，不在压缩包内重复打包。
\end{itemize}

\subsection{提交结果工作簿}

\begin{table}[H]
  \centering
  \caption{提交结果工作簿、工作表与生成入口}
  \label{tab:app-workbooks}
  \small
  \setlength{\tabcolsep}{4pt}
  \begin{tabularx}{\textwidth}{l>{\raggedright\arraybackslash}Xl}
    \toprule
    文件 & 工作表 & 生成入口 \\
    \midrule
    \texttt{results/result1.xlsx} & 计划购电量、充放电量 & \texttt{scripts/q1\_solve.py} \\
    \texttt{results/result2.xlsx} & 计划购电量、充放电量、紧急购电量 & \texttt{scripts/q2\_finalize\_E.py} \\
    \texttt{问题三/result3.xlsx} & 计划购电量、调整购电量、充放电量、紧急购电量 & \texttt{问题三/代码/} \\
    \texttt{问题四/result4-2.xlsx} & 计划购电量、充放电量、紧急购电量 & 问题四链路 \\
    \texttt{问题四/result4-3.xlsx} & 计划购电量、调整购电量、充放电量、紧急购电量 & 问题四链路 \\
    \bottomrule
  \end{tabularx}
\end{table}

五个工作簿的工作表名与官方模板逐一对应，其中前四个已按模板行序核对；问题三的 \texttt{result3.xlsx} 由问题三链路生成，未纳入本文的核验清单（见 \S\ref{app:audit}）。

\subsection{计算脚本}

以下脚本按问题分组；\textbf{每问只有一个复现入口}，其余脚本为该问的分析、配图或档案工具。

\begin{itemize}
  \sloppy \footnotesize
  \item \textbf{问题一}：主程序 \texttt{scripts/q1\_solve.py} 是本问的唯一入口，完成线性规划装配、求解、独立校验并写出 \texttt{result1.xlsx}；敏感性分析由 \texttt{scripts/q1\_analysis.py} 完成，涵盖无储能基准、容量与功率敏感性以及边际价值；配图由 \texttt{scripts/q1\_figures.py} 与 \texttt{scripts/q1\_analysis\_figures.py} 生成，前者绘制代表日输入与最优策略，后者绘制敏感性与边际价值曲线。
  \item \textbf{问题二入口}：主程序 \texttt{scripts/q2\_finalize\_E.py} 是本问的唯一入口，复现评价期总费用、生成三张指定日期结果表并写出 \texttt{result2.xlsx}。
  \item \textbf{问题二辅助}：选参由 \texttt{scripts/q2\_dow\_N\_joint.py} 完成，回放 36 个候选并给出逐月 walk-forward 日程；月度费用与风险档案由 \texttt{scripts/q2\_V\_dow\_N\_profile.py} 生成；\texttt{scripts/q2\_model.py} 回放早期的对照方案 A--D，\texttt{scripts/q2\_improve.py} 则用于预测器形式与执行规则的改进对照。
  \item \textbf{问题三}：\texttt{问题三/代码/} 下的模块实现附件 3 的预报通道、分档收缩回归、调整层线性规划与分段执行；\texttt{问题三/图片/} 存放该问的结果图。
  \item \textbf{问题四}：\texttt{问题四/代码/} 下的模块实现 48 小时滚动时域的计划与调整线性规划、尽限执行与全年回放；\texttt{scripts/q4\_export\_pt.py} 导出问题二预测器的中心预测序列，作为该问链路的输入。
  \item \textbf{图源文件}：\texttt{问题一/图片/fig*.fig} 与 \texttt{问题二/图片/fig*.fig} 为问题一、问题二配图的图源文件。
\end{itemize}

问题四的运行需先执行 \texttt{q4\_export\_pt.py} 生成需求代理序列，该依赖为单向前置关系；除此之外，四问的复现入口互不重叠，不存在同一问被两个脚本各自声称唯一入口的情况。

\subsection{中间结果}

\begin{itemize}
  \sloppy \small
  \item 问题一：\texttt{results/q1\_summary.json}（解与校验）、\texttt{results/q1\_analysis.json}（敏感性与边际价值）、\texttt{results/q1\_solution.npz}（求解器输出）。
  \item 问题二：\texttt{results/q2\_dow\_N\_joint.json}（选参日程）、\texttt{results/q2\_V\_dow\_N\_profile.json}（月度与风险档案）、\texttt{results/q2\_E\_tables.json}（指定日期表格）。
  \item 问题四的输入之一：\texttt{results/q4\_v3/F\_pt.csv}，即问题二预测器的中心预测序列（$365\times144$）。
\end{itemize}

上述中间结果均由其对应脚本自动生成，不含手工编辑的数值。除 \texttt{F\_pt.csv} 是问题四链路的必要输入外，其余中间结果不参与任何论文数值的最终生成：论文中的数字一律由脚本在运行时直接输出。

\section{软件环境与依赖版本}

\begin{table}[H]
  \centering
  \caption{运行环境与依赖版本（实际使用版本）}
  \label{tab:app-env}
  \small
  \begin{tabularx}{\textwidth}{l>{\raggedright\arraybackslash}X}
    \toprule
    组件 & 版本／说明 \\
    \midrule
    Python & 3.13.14 \\
    NumPy & 2.5.2 \\
    SciPy & 1.18.1（线性规划经 \texttt{scipy.optimize.linprog} 调用 HiGHS 求解器） \\
    pandas & 3.0.5（读取 \texttt{.xlsx} 输入） \\
    openpyxl & 3.1.5（写出提交工作簿） \\
    MATLAB & 用于问题一、问题二配图的图源文件（\texttt{.fig}），不作为求解器使用 \\
    \bottomrule
  \end{tabularx}
\end{table}

全部脚本以项目根目录为基准解析相对路径（即以脚本文件所在目录的两级父目录为根），源码中不含任何个人绝对路径，可在任意目录放置后直接运行。

\section{论文图表与数据、脚本的对应关系}

\begin{table}[H]
  \centering
  \caption{论文主要图表与配置、数据文件、生成脚本的对应关系}
  \label{tab:app-map}
  \footnotesize
  \setlength{\tabcolsep}{4pt}
  \begin{tabularx}{\textwidth}{l>{\raggedright\arraybackslash}Xl}
    \toprule
    论文图表 & 对应配置与数据 & 生成脚本 \\
    \midrule
    表~\ref{tab:q1-t1}、表~\ref{tab:q1-t2} & 问题一最优解；\texttt{data/附件1.xlsx} & \texttt{q1\_solve.py} \\
    表~\ref{tab:q1-compare} & 无储能基准与主方案 & \texttt{q1\_analysis.py} \\
    图~\ref{fig:q1-input}--图~\ref{fig:q1-sens}（问题一） & 代表日输入、最优策略、敏感性与边际价值 & \texttt{q1\_figures.py}、\texttt{q1\_analysis\_figures.py} \\
    表~\ref{tab:q2-sched} & 36 候选逐月 walk-forward & \texttt{q2\_dow\_N\_joint.py} \\
    表~\ref{tab:q2-total}、表~\ref{tab:q2-t1}--表~\ref{tab:q2-t3} & 主方案 E；\texttt{data/附件1,2.xlsx} & \texttt{q2\_finalize\_E.py} \\
    表~\ref{tab:q2-month}、表~\ref{tab:q2-risk} & 主方案月度与风险档案 & \texttt{q2\_V\_dow\_N\_profile.py} \\
    表~\ref{tab:q2-ladder}、图~\ref{fig:q2-ladder} & 对照方案 A--E 的水平比较 & \texttt{q2\_model.py}（回放）、\texttt{q\_paper\_figures.py}（绘图） \\
    图~\ref{fig:q2-frame}--图~\ref{fig:q2-emerg}（问题二） & 主方案 E 的框架、预测、指定日、全年与风险 & \texttt{问题二/图片/fig*.fig} 导出 \\
    表~\ref{tab:q3-matrix}、表~\ref{tab:q3-b3} & 问题三 $2\times2$ 矩阵与 B3 费用分解 & \texttt{问题三/代码/} \\
    图~\ref{fig:q3-matrix}--图~\ref{fig:q3-month}（问题三） & 矩阵、MAE 与费用瀑布、日内形态、门控与逐月 & \texttt{问题三/图片/fig*.pdf} \\
    表~\ref{tab:q4-struct}、图~\ref{fig:q4-struct} & 附件 4 价格结构（可直接复算） & 复算脚本 \\
    表~\ref{tab:q4-main}、表~\ref{tab:q4-horizon} & 主方案 F\_48／G\_48 与时域对照 & 问题四链路 \\
    图~\ref{fig:q4-horizon}、图~\ref{fig:q4-ladder}、图~\ref{fig:q4-regret} & 时域机制、方案阶梯、参照对比 & \texttt{问题四/照片/fig*.pdf} \\
    \bottomrule
  \end{tabularx}
\end{table}

\section{校验与审计的分项说明}
\label{app:audit}

\begin{table}[H]
  \centering
  \caption{审计分项、覆盖范围与结果}
  \label{tab:app-audit}
  \footnotesize
  \setlength{\tabcolsep}{4pt}
  \begin{tabularx}{\textwidth}{l>{\raggedright\arraybackslash}Xl}
    \toprule
    审计项 & 覆盖范围 & 结果 \\
    \midrule
    物理守恒（区间能量平衡） & 问题一单日；问题二评价期 334 天；问题四全部 20 条轨迹 & 最大残差 $\le1.2\times10^{-13}$ kWh（Q2）、$<10^{-7}$（Q4） \\
    储能电量边界 & 同上 & 全程位于 $[1200,10800]$ kWh \\
    充放电功率上限 & 同上 & 无越界；同时充放电重叠量为 0 \\
    费用账本 & 问题一全天费用；问题二计划费与紧急费；问题四两分支合同费与附加费 & 与独立复算路径逐位吻合 \\
    时间映射 & \texttt{result1.xlsx}、\texttt{result2.xlsx} 与字面时钟区间 & 模板行标签整体后移一个区间，两口径相差 10 分钟 \\
    结果模板 & 四个已生成工作簿的工作表名与行序 & 与官方模板一致；\texttt{result3} 未纳入本次核验 \\
    可复现性 & 问题一、问题二主链路 & 可由原始附件直接复现；问题四需先运行 \texttt{q4\_export\_pt.py} \\
    \bottomrule
  \end{tabularx}
\end{table}

需要区分``一致''与``正确''：``沿用旧脚本得到逐位相同的数值''只能证明两次运行一致，不能证明旧口径本身正确。因此上表中凡涉及账本与模板的项，均由\emph{不调用调度代码}的独立路径重算，而不是与原脚本自比。

\section{核心程序}

本节给出问题一、问题二与问题四核心模块的\textbf{完整源码}（三个文件即对应链路的核心实现，可独立阅读），其余辅助模块（问题一敏感性分析、问题二选参与风险档案、问题三预报融合与调整层、问题四全量回放与核验）见上文 \S 的脚本清单，随支撑材料按目录结构提交。

\subsection{问题一：线性规划装配、求解与独立复核}
\label{app:q1-code}

\VerbatimInput[frame=single,numbers=left,fontsize=\scriptsize,label=\texttt{code/q1\_core.py}]{code/q1_core.py}

\subsection{问题二：预测—计划—执行三层框架核心代码}
\label{app:q2-code}

\VerbatimInput[frame=single,numbers=left,fontsize=\scriptsize,label=\texttt{code/q2\_core.py}]{code/q2_core.py}

\subsection{问题四：48 小时滚动时域计划、调整与执行核心代码}
\label{app:q4-code}

\VerbatimInput[frame=single,numbers=left,fontsize=\scriptsize,label=\texttt{code/q4\_core.py}]{code/q4_core.py}
